\section{Optical Printing and Print Stock}
\label{sec:print}

\subsection{Why the Print Stage Is Not Optional}

The negative is a low-contrast, orange-masked, inverted record with a gamma near
$0.6$. Nothing about it resembles a photograph. Every visual quality that people
mean when they say ``film look'' --- the contrast, the highlight roll-off, the
specific way saturated colors desaturate as they approach white, the black point
--- is created in the \emph{second} transfer, when the negative is projected onto
print stock.

Simulations that stop at the negative and then apply a contrast curve produce a
result that is qualitatively wrong in a way that is easy to see and hard to
name. The reason is that the print transfer is not a curve applied to the
negative's \emph{output}; it is a curve applied to the negative's
\emph{transmittance}, mediated by a crosstalk matrix, with the negative's three
records interacting. That composition is not expressible as three independent
one-dimensional curves.

\subsection{Negative Transmittance and Printing Density}

The negative's transmittance in record $k$ is, by definition of density,

\begin{equation}
T_k = 10^{-D_k}.
\end{equation}

The print stock's red-sensitive layer, however, does not see $T_R$. It sees the
integral of the enlarger's spectral power distribution, the negative's total
spectral transmittance (all three dyes overlap), and the layer's own spectral
sensitivity. Reducing that integral to three channels yields a
\emph{printing density} matrix $\mathbf{C}$ that maps the negative's measured
densities to the effective densities seen by each print layer:

\begin{equation}
D^{\mathrm{eff}}_c \;=\; \sum_{k} C_{ck}\, D_k .
\label{eq:printdensity}
\end{equation}

The off-diagonal terms of $\mathbf{C}$ are the \emph{unwanted absorptions}: the
cyan dye absorbs some green, the magenta dye absorbs some blue, and so on. For a
modern color negative printed on a modern print stock,

\begin{equation}
\mathbf{C} \approx
\begin{pmatrix}
1.000 & 0.086 & 0.028 \\
0.041 & 1.000 & 0.113 \\
0.017 & 0.052 & 1.000
\end{pmatrix}.
\label{eq:cmatrix}
\end{equation}

$\mathbf{C}$ is derived offline from spectral dye absorption data
\cite{kodak2383,james1977} and print
stock spectral sensitivities; it is the principal artifact of the reduced
spectral strategy announced in Section~\ref{sec:background}.

The print exposure is then

\begin{equation}
\logE'_c \;=\; \log_{10} L_c \;-\; D^{\mathrm{eff}}_c \;+\; \log_{10} t_{\mathrm{print}},
\label{eq:printexposure}
\end{equation}

where $L_c$ is the enlarger channel intensity (set by the printer lights of
Section~\ref{sec:printerlights}) and $t_{\mathrm{print}}$ is the print exposure
time, which is the overall density control.

\subsection{Automatic Mask Removal}

Observe that the negative's $\Dmin$ enters \eqref{eq:printdensity} as a constant
vector, hence enters \eqref{eq:printexposure} as a constant offset in
$\logE'$. It is therefore exactly compensable by a constant adjustment of $L_c$
--- which is precisely what a laboratory does when it establishes the printer
light values for a given stock. The orange mask is not removed by a special
operation; it is removed by \emph{balancing}, and the model reproduces that
automatically.

\EM{} computes the neutral printer lights $L_c^{\mathrm{aim}}$ at profile-load
time by requiring that a scene-referred neutral at $E = 0.18$ reproduce the
print stock's aim density $\mathbf{D}'_{\mathrm{aim}}$:

\begin{equation}
\log_{10} L_c^{\mathrm{aim}} \;=\; D'^{-1}_{\mathrm{print},c}\!\left(D'_{\mathrm{aim},c}\right) \;+\; D^{\mathrm{eff}}_c\big|_{E=0.18},
\label{eq:aimbalance}
\end{equation}

where $D'^{-1}_{\mathrm{print},c}$ is the numerical inverse of the print stock's
characteristic curve, obtained by four Newton iterations from the straight-line
estimate. This calculation is performed once per (negative stock, print stock)
pair and cached, and is what guarantees AC-5 across all $10 \times 4$ shipping
combinations without hand-tuning each one.

\subsection{The Print Characteristic Curve}

The print stock uses the identical formulation of \eqref{eq:charcurve}, with
substantially different parameters:

\begin{equation}
D'_c = D'_{\min,c} + \softp{\kappa'_{t,c}}{u'_c} - \softp{\kappa'_{s,c}}{u'_c - \Delta D'_c},
\label{eq:printcurve}
\end{equation}

with $u'_c = \gamma'_c(\logE'_c - x'_{0,c})$. Typical print stock parameters
are $\gamma' \in [2.4, 3.2]$, $\Delta D' \in [2.2, 2.6]$, and a
\emph{much} harder shoulder than the negative, $\kappa'_s \approx 0.06$, which
is what produces the abrupt, characteristic approach to maximum black.

Because the print inverts, the negative's shoulder (scene highlights, high
negative density, low print exposure) maps to the print's \emph{toe}. The
graceful highlight roll-off that film is celebrated for is therefore the
composition of two soft regions --- the negative's shoulder and the print's toe
--- and reproducing it requires both. Applications that model only one produce
highlights that roll off half as gracefully, which reads as "digital."

The overall system gamma is the product of the straight-line gammas modulated by
the crosstalk matrix; for the representative pairing,

\begin{equation}
\gamma_{\mathrm{sys}} = \gamma_{\mathrm{neg}} \cdot \gamma_{\mathrm{print}} \approx 0.61 \times 2.80 \approx 1.71,
\end{equation}

consistent with the $1.5$--$1.8$ range appropriate for a dark-surround viewing
condition \cite{hunt2004}.

\begin{figure}[htbp]
\centering
\begin{tikzpicture}
\begin{axis}[
  width=\columnwidth, height=5.2cm,
  xlabel={scene $\logE$ (relative)},
  ylabel={print density $D'$},
  xmin=-2.6, xmax=1.4, ymin=0, ymax=2.8,
  grid=both, grid style={draw=emLine!20},
  legend style={font=\scriptsize, at={(0.98,0.98)}, anchor=north east, draw=emLine!40},
  label style={font=\scriptsize}, tick label style={font=\scriptsize},
  samples=250,
]
% Composite: print curve driven by (aim - gamma_neg*(x - x0n)) as effective print logE
% Effective print logE' = k - gamma_neg*(x + 2.10)   (negative density rises with x)
\addplot[emAcc, very thick, domain=-2.6:1.4]
  {0.06 + 0.07*ln(1+exp((2.80*((1.05 - 0.61*(x+2.10)) + 0.55))/0.07))
        - 0.06*ln(1+exp((2.80*((1.05 - 0.61*(x+2.10)) + 0.55) - 2.45)/0.06))};
\addlegendentry{print density $D'$}
\end{axis}
\end{tikzpicture}
\caption{Composite scene-to-print transfer, obtained by driving \eqref{eq:printcurve} with the negative density of \eqref{eq:charcurve}. The rising left portion is scene shadow rendered as print black; the falling right portion is scene highlight rendered as print white. The soft knee at the right is the composition of the negative's shoulder with the print's toe.}
\label{fig:printcomposite}
\end{figure}

\subsection{Print Stock Profiles}

\begin{table}[htbp]
\caption{Shipping Print Stock Profiles}
\label{tab:printstocks}
\begin{center}
\renewcommand{\arraystretch}{1.15}
\begin{tabularx}{\columnwidth}{@{}l>{\raggedright\arraybackslash}X@{}}
\toprule
\textbf{Profile} & \textbf{Character} \\
\midrule
Standard 2383-type & The default. $\gamma' = 2.80$, moderate saturation, warm-leaning neutral axis, gentle toe. The reference against which everything else is judged. \\
Premium 2393-type & $\gamma' = 3.05$, higher saturation via steeper cross-terms, deeper maximum density, harder shoulder. More contrast, less latitude. \\
Fuji 3513-type & $\gamma' = 2.72$, cooler neutral axis, distinctly different green and cyan rendering, softer approach to maximum density. \\
Fuji 3521-type & $\gamma' = 2.90$, higher saturation Fuji rendering, characteristic magenta handling. \\
Bypass (scan look) & No print transfer. Negative density is inverted and normalized only. Produces the flat, low-contrast look of an unadjusted lab scan --- included deliberately, because it is what many users have actually seen and expect. \\
\bottomrule
\end{tabularx}
\end{center}
\end{table}

\subsection{Silver Retention (Bleach Bypass and ENR)}

Skipping or shortening the bleach step leaves metallic silver in the print
alongside the dye image. Silver is spectrally neutral, so the result is a
neutral density image added to the color image: contrast rises sharply,
saturation falls, and blacks become very dense.

The model is direct. Let $\bar{D}'$ be the neutral (silver) density, which is
proportional to total development:

\begin{equation}
\bar{D}' = \frac{1}{3}\sum_c \frac{D'_c - D'_{\min,c}}{\Delta D'_c} \cdot \Delta D'_{\mathrm{Ag}},
\end{equation}

and the retained-silver print density is

\begin{equation}
D'^{\mathrm{Ag}}_c = D'_c + \varrho\, \bar{D}',
\label{eq:silverretention}
\end{equation}

with $\varrho \in [0, 1]$ the retention fraction. $\varrho = 1$ is full bleach
bypass; intermediate values correspond to the ENR and ACE processes, which
retain a controlled fraction. This single scalar reproduces the entire family
of silver-retention looks, and it reproduces their characteristic property ---
that desaturation is strongest in the shadows, where $\bar{D}'$ is largest
relative to the color record --- automatically.

\subsection{Print Density to Display}

The print is viewed by transmission (projection) or reflection. In either case
the luminance reaching the eye is proportional to $10^{-D'}$, so the display
transform is

\begin{equation}
Y_c = \frac{10^{-D'_c} - 10^{-D'_{\max,c}}}{10^{-D'_{\min,c}} - 10^{-D'_{\max,c}}},
\label{eq:printtodisplay}
\end{equation}

which normalizes the print's own $\Dmin$ to display white and its $\Dmax$ to
display black. The subtraction of the $\Dmax$ term is important: without it, the
print's finite maximum density leaves a lifted black that compounds with the
display's own black level and looks washed out on an OLED panel.

The resulting linear $Y$ is converted from the print stock's primaries to
Display P3 and encoded. A viewing-condition parameter applies an optional
surround compensation exponent \cite{fairchild2013} (default $1.0$ for a phone
screen viewed in
normal room light, $0.9$ available for a simulated dark-surround projection
look).
